\documentclass[11pt]{amsart}

\usepackage[T1]{fontenc}
\usepackage{lmodern}
\usepackage{microtype}
\usepackage{amsmath,amssymb,amsthm}
\usepackage[hidelinks]{hyperref}

\hypersetup{
  pdftitle={Two Residual Cases of Simmons's Polygonal Bigraph Conjecture}
}

\newtheorem{theorem}{Theorem}
\newtheorem{lemma}[theorem]{Lemma}
\newtheorem{corollary}[theorem]{Corollary}

\newcommand{\Z}{\mathbb Z}

\title[Two residual polygonal bigraphs]
{Two Residual Cases of Simmons's Polygonal Bigraph Conjecture}

\date{}

\subjclass[2020]{05C45}
\keywords{
Hamilton cycle,
2-edge-Hamiltonian graph,
cubic bipartite graph,
polygonal bigraph
}

\begin{document}

\begin{abstract}
Simmons conjectured that every polygonal bigraph $P_{m,k}$ is
2-edge-Hamiltonian. His addendum left the isomorphic pair
$P_{34,9}$ and $P_{34,14}$ unresolved, while asserting the conjecture for
all polygonal bigraphs on at most $66$ vertices. Its residual table also
omitted $P_{32,14}$, the only omitted entry not removed by the stated
isomorphism reductions; as $P_{32,14}$ has $64$ vertices, the $66$-vertex
assertion was not in fact established. We give five Hamilton cycles for
each of $P_{34,9}$ and $P_{32,14}$ whose edge-pair coverage proves that
both graphs are 2-edge-Hamiltonian. An explicit isomorphism settles
$P_{34,14}$. Consequently the conjecture holds for every polygonal
bigraph on at most $68$ vertices, which both repairs the $66$-vertex
claim and extends it. For $P_{32,14}$ the conclusion is not new: it is
the case $s=4$ of the family $P_{8s,4s-2}$, treated in a companion note
in this collection, and the certificate given here is an independent one.
What is new is the class $P_{34,9}\cong P_{34,14}$ together with the
finite audit. This is a finite verification, not a proof for arbitrary $m$.
\end{abstract}

\maketitle

\section{Statement}

A graph is \emph{2-edge-Hamiltonian} if every two distinct edges lie on a
common Hamilton cycle. For $m\ge 3$ and
$1\le k\le \lfloor(m-1)/2\rfloor$, let $P_{m,k}$ have bipartition
\[
 X=\{x_i:i\in\Z_m\},
 \qquad
 Y=\{y_i:i\in\Z_m\},
\]
and edges, with subscripts modulo $m$,
\[
 a_i=x_i y_i,
 \qquad
 b_i=x_i y_{i-1},
 \qquad
 c_i=x_i y_{i+k}.
\]
Thus $P_{m,k}$ is the usual $2m$-cycle together with the chords of offset
$2k+1$.

Simmons conjectured that every $P_{m,k}$ is 2-edge-Hamiltonian
\cite{Simmons2020}. His addendum identified
$P_{34,9}\cong P_{34,14}$ as an unresolved isomorphism class
\cite{SimmonsAddendum}; the isomorphism reductions used there originate in
\cite{Simmons2016}. The same addendum asserted the result through
$66$ vertices. Its displayed residual table, however, omitted the entries
$P_{31,11}$, $P_{32,8}$, and $P_{32,14}$. The first two reduce to settled
parameters, but $P_{32,14}$ does not, so the assertion for $66$ vertices
was not in fact established: $P_{32,14}$ has $64$ vertices.
Theorem~\ref{thm:main} therefore both repairs the $66$-vertex claim and
extends it to $68$.

For $P_{32,14}$ the conclusion of Theorem~\ref{thm:main} is not new. That
graph is the case $s=4$ of the family $P_{8s,4s-2}$, every member of which
is proved 2-edge-Hamiltonian in the companion note \cite{Companion} in
this collection; the certificate we give below is an independent one, and
we include it so that the $66$-vertex range is repaired within the present
note. What is new here is the isomorphism class
$P_{34,9}\cong P_{34,14}$, which lies in no family treated in
\cite{Companion}, together with the finite audit of the residual
parameters with $m\le34$.

\begin{theorem}\label{thm:main}
The polygonal bigraphs
\[
 P_{32,14},
 \qquad
 P_{34,9},
 \qquad
 P_{34,14}
\]
are 2-edge-Hamiltonian.
\end{theorem}

\begin{corollary}\label{cor:68}
Every polygonal bigraph on at most $68$ vertices is 2-edge-Hamiltonian.
\end{corollary}

\section{A covering certificate}

For $S\subseteq\Z_m$ and $t\in\Z_m$, write
\[
 S+t=\{s+t:s\in S\}.
\]
If
\[
 \Z_m=A\mathbin{\dot\cup}B\mathbin{\dot\cup}C,
\]
define
\[
 M(A,B,C)
 =
 \{a_i:i\in A\}
 \cup
 \{b_i:i\in B\}
 \cup
 \{c_i:i\in C\}.
\]
This set meets every vertex of $X$ once. Its $Y$-endpoints are indexed by
$A$, $B-1$, and $C+k$; hence it is a perfect matching precisely when those
three sets partition $\Z_m$.

\begin{lemma}\label{lem:cover}
Let $M_0,\dots,M_4$ be perfect matchings of a cubic graph $G$. Suppose each
\[
 H_r=G-M_r
\]
is a Hamilton cycle and, for every two distinct edges $e,f$, some $M_r$
contains neither $e$ nor $f$. Then $G$ is 2-edge-Hamiltonian.
\end{lemma}

\begin{proof}
For that index $r$, both edges belong to the Hamilton cycle $H_r$.
\end{proof}

\section{The two certificates}

\subsection{The graph \texorpdfstring{$P_{34,9}$}{P34,9}}

For $r=0,\dots,4$, take the following subsets of $\Z_{34}$ and set
\[
 A_r=\Z_{34}\setminus(B_r\cup C_r):
\]
\[
\begin{array}{c|c|c}
 r & B_r & C_r \\
\hline
 0 & \varnothing & \Z_{34} \\
 1 & \Z_{34} & \varnothing \\
 2 & \{25,26\} & \{8,17,24,33\} \\
 3 & \{16,24\} & \{7,15,23,32\} \\
 4 & \{15,23\} & \{6,14,22,31\}
\end{array}
\tag{3.1}\label{tab:349}
\]
Put
\[
 M_r=M(A_r,B_r,C_r),
 \qquad
 H_r=P_{34,9}-M_r.
\]
Directly from \eqref{tab:349}, the sets $A_r$, $B_r-1$, and $C_r+9$
partition $\Z_{34}$, so every $M_r$ is a perfect matching.

The cycle $H_0$ is the outer $68$-cycle. The cycle $H_1$ consists of all
$a$- and $c$-edges; its successive $X$-indices differ by $-9$, and
\[
 \gcd(34,9)=1.
\]
Thus $H_0$ and $H_1$ are Hamilton cycles. The three traces in
Appendix~\ref{app:traces} show that $H_2,H_3,H_4$ are also Hamilton
cycles.

Now $H_0$ contains every $a$- and $b$-edge, while $H_1$ contains every
$a$- and $c$-edge. Only pairs $\{b_i,c_j\}$ remain. The sets
$B_2,B_3,B_4$ are pairwise disjoint, as are $C_2,C_3,C_4$. Hence a fixed
$b_i$ is absent from at most one of $H_2,H_3,H_4$, and the same holds for
a fixed $c_j$. At least one of these three cycles therefore contains both
edges. Lemma~\ref{lem:cover} proves that $P_{34,9}$ is
2-edge-Hamiltonian.

\subsection{The graph \texorpdfstring{$P_{32,14}$}{P32,14}}

For $r=0,\dots,4$, take the following subsets of $\Z_{32}$ and set
\[
 B_r=\Z_{32}\setminus(A_r\cup C_r):
\]
\[
\begin{array}{c|c|c}
 r & A_r & C_r \\
\hline
 0 & \varnothing & \Z_{32} \\
 1 & \Z_{32} & \varnothing \\
 2 & \{0,16\} & \{1,17\} \\
 3 & \{15,31\} & \{0,16\} \\
 4 & \{14,30\} & \{15,31\}
\end{array}
\tag{3.2}\label{tab:3214}
\]
Put
\[
 M_r=M(A_r,B_r,C_r),
 \qquad
 H_r=P_{32,14}-M_r.
\]
For every row, $A_r$, $B_r-1$, and $C_r+14$ partition $\Z_{32}$, so every
$M_r$ is a perfect matching.

The cycle $H_0$ is the outer $64$-cycle. The cycle $H_1$ consists of all
$b$- and $c$-edges; up to orientation, its successive $X$-indices differ
by $15$, and
\[
 \gcd(32,15)=1.
\]
The remaining traces in Appendix~\ref{app:traces} show that
$H_2,H_3,H_4$ are Hamilton cycles.

Here $H_0$ contains every $a$- and $b$-edge, while $H_1$ contains every
$b$- and $c$-edge. Only pairs $\{a_i,c_j\}$ remain. The sets
$A_2,A_3,A_4$ are pairwise disjoint, as are $C_2,C_3,C_4$. The same
three-cycle argument proves that every such pair lies on a Hamilton cycle.
Thus $P_{32,14}$ is 2-edge-Hamiltonian.

\section{Isomorphism and the finite range}

Let primed symbols denote $P_{34,14}$. Define
\[
 \phi(x_i)=x'_{19i},
 \qquad
 \phi(y_i)=y'_{19i-1}.
\]
Since $19$ is a unit modulo $34$, this is a bijection, and
\[
 \phi(a_i)=b'_{19i},
 \qquad
 \phi(b_i)=c'_{19i},
 \qquad
 \phi(c_i)=a'_{19i}.
\]
Indeed,
\[
 19(i-1)-1\equiv19i+14,
 \qquad
 19(i+9)-1\equiv19i
 \pmod{34}.
\]
Thus
\[
 P_{34,9}\cong P_{34,14},
\]
completing the proof of Theorem~\ref{thm:main}.

For completeness, we record the finite audit needed for
Corollary~\ref{cor:68}. Put
\[
 S_k=\{0,-1,k\}\subseteq\Z_m.
\]
If $u\in\Z_m^\times$ and $d\in\Z_m$ satisfy
\[
 uS_k+d=S_\ell,
\]
then
\[
 x_i\longmapsto x'_{ui},
 \qquad
 y_j\longmapsto y'_{uj+d}
\]
is an isomorphism
\[
 P_{m,k}\cong P_{m,\ell}.
\]

After Simmons's general results, the residual arithmetic condition is
\[
 2<\gcd(m+1,k+1)<k+1.
\]
This reduction is recorded in \cite{Simmons2020,SimmonsAddendum}.
For $m\le34$, its solutions and their
reductions are exactly the following. An arrow
\[
 k\to\ell\;(u,d)
\]
records $uS_k+d=S_\ell$.

\begin{center}
\small
\begin{tabular}{c|l}
$m$ & residual parameters and reductions \\
\hline
14 & $5\to3\;(11,0)$ \\
19 & $7$ (settled in \cite{Simmons2020}) \\
20 & $5,8$ (settled in \cite{Simmons2020}) \\
23 & $8\to3\;(20,0)$ \\
24 & $9$ (settled in \cite{Simmons2020}) \\
26 & $5\to4\;(21,-1)$, $11\to6\;(19,-1)$ \\
27 & $7\to4\;(23,0)$, $11\to5\;(22,0)$ \\
29 & $8\to12\;(13,12)$, $11\to12\;(12,12)$ \\
31 & $11\to12\;(13,12)$ \\
32 & $5\to13\;(19,0)$, $8\to7\;(7,7)$,
     $11\to3\;(29,0)$, $14$ \\
34 & $13\to12\;(21,-1)$, $9\to14\;(19,-1)$
\end{tabular}
\end{center}

Every arrow except $9\to14$ lands in a parameter covered by the results of
\cite{Simmons2020}. The unreduced entry $P_{32,14}$ and the class
\[
 P_{34,9}\cong P_{34,14}
\]
are precisely those proved in Theorem~\ref{thm:main}. Hence every
$P_{m,k}$ with $m\le34$ is 2-edge-Hamiltonian. Since $P_{m,k}$ has $2m$
vertices, Corollary~\ref{cor:68} follows.

\appendix

\section{Hamilton traces}\label{app:traces}

Each row below is a cyclic list of successive $X$-indices in the indicated
spanning $2$-factor. Consecutive entries, including the last and first,
have a common $Y$-neighbor through two retained edges. Since each row
contains every $X$-index exactly once, it certifies a Hamilton cycle.

\subsection*{\texorpdfstring{$P_{34,9}$}{P34,9}}

{\scriptsize
\begin{align*}
T_{34,2}={}&
(0,10,20,30,6,16,25,1,11,21,31,7,17,18,28,4,14,\\
&24,15,5,29,19,9,8,32,22,12,2,26,27,3,13,23,33),
\\[2pt]
T_{34,3}={}&
(0,10,20,30,6,15,5,29,19,9,33,32,22,12,2,26,16,\\
&17,27,3,13,23,14,4,28,18,8,7,31,21,11,1,25,24),
\\[2pt]
T_{34,4}={}&
(0,10,20,30,6,7,17,27,3,13,22,12,2,26,16,15,25,\\
&1,11,21,31,32,8,18,28,4,14,5,29,19,9,33,23,24).
\end{align*}
}

\subsection*{\texorpdfstring{$P_{32,14}$}{P32,14}}

{\scriptsize
\begin{align*}
T_{32,2}={}&
(0,14,28,10,24,6,20,2,17,3,21,7,25,11,29,15,\\
&16,30,12,26,8,22,4,18,1,19,5,23,9,27,13,31),
\\[2pt]
T_{32,3}={}&
(0,18,4,22,8,26,12,30,31,13,27,9,23,5,19,1,\\
&16,2,20,6,24,10,28,14,15,29,11,25,7,21,3,17),
\\[2pt]
T_{32,4}={}&
(0,18,4,22,8,26,12,30,29,11,25,7,21,3,17,31,\\
&16,2,20,6,24,10,28,14,13,27,9,23,5,19,1,15).
\end{align*}
}

\begin{thebibliography}{9}

\bibitem{Simmons2016}
G.~J. Simmons,
\emph{Hamilton-laceable properties of polygonal bigraphs},
Congr. Numer. \textbf{226} (2016), 121--138.

\bibitem{Simmons2020}
G.~J. Simmons,
\emph{Two related (?) 2-edge-Hamiltonian bigraph conjectures},
Bull. Inst. Combin. Appl. \textbf{88} (2020), 98--117.
\url{https://bica.the-ica.org/Volumes/88/Reprints/BICA2019-31-Reprint.pdf}

\bibitem{SimmonsAddendum}
G.~J. Simmons,
\emph{Two related (?) 2-edge-Hamiltonian bigraph conjectures addendum},
Bull. Inst. Combin. Appl. \textbf{89} (2020), 102--103.
\url{https://bica.the-ica.org/Volumes/89/Reprints/Gus-Reprint.pdf}

\bibitem{Companion}
\emph{The polygonal bigraphs $P_{8s,4s-2}$ are 2-edge-Hamiltonian},
companion note in this collection, file
\texttt{p8s-4sminus2-\allowbreak confirms-simmons-hamiltonian-conjecture}
in \texttt{notes/}.

\end{thebibliography}

\end{document}